\paragraph{Data and splits.} For each seed $s\in\{0,\dots,4\}$ we generate \rhval{const/n_train_seq} distinct provable training sequents with \rhval{const/train_lo}--\rhval{const/train_hi} connectives and \rhval{const/n_val_seq} further held-out sequents of the same range for validation (top-1 accuracy of the policy against the optimal set, and the \emph{val} task). The test tasks contain \rhval{const/n_probs} problems each, drawn from RNGs seeded by (task, seed) and thus identical across systems: \emph{test\_11\_20}, \emph{test\_21\_40} and \emph{test\_41\_70} (connective ranges 11--20, 21--40, 41--70, all strictly larger than any training sequent), plus \emph{val\_3\_10} (training range, held out) and \emph{tune\_21\_35}, used only to choose the heuristic variant. Changing the seed changes the training set, the test problems and the model initialisation.

\paragraph{Metrics.} \emph{solved}: fraction of problems proved within \rhval{const/budget} expansions (primary, higher is better); \emph{mean\_exp}: mean number of expansions with failures counted as the budget (lower is better); \emph{val\_top1} and \emph{val\_loss}: policy accuracy and loss on held-out training-range states. All systems use the same engine, budget, problems and metric definition.

\paragraph{Training and tuning.} Both policies: AdamW (peak learning rate \rhval{const/lr}, weight decay $0.01$, one-cycle schedule), batch \rhval{const/batch}, \rhval{const/train_steps} steps, gradient clipping at 1; no hyperparameter of either policy was tuned, and none was ever selected on a test bin. The heuristic was tuned: three ordering variants were run on \emph{tune\_21\_35} with all five seeds (Table~\ref{tab:tune}), and the best (\emph{large}) was used everywhere; this tuning bin is in the size regime of the test bins while the policies saw only training-range data, so the baseline has the advantage of access to the target size regime. Hardware: shared CPU, \rhval{const/threads} threads per process, at most two processes at once. A policy is trained once per (system, seed) and its checkpoint is reused by all evaluation runs.

\paragraph{Pre-registered hypotheses and tests.} (H1) transformer-guided search solves more problems than BFS in each test bin; (H2) it solves more than the heuristic at 41--70; (H3) the transformer-minus-heuristic gap decreases monotonically from the 11--20 bin to the 41--70 bin; (H4) at 41--70 best-first search with the policy beats greedy decoding of the same policy. The tests are paired $t$-tests over the five seeds (the problems of a seed are identical for all systems), produced by \texttt{rh compare} and by \texttt{method/analyze.py}. With five seeds these tests have low power for small effects and the $p$-values are not corrected for multiple comparisons.

\begin{table}[t]\centering\caption{Hand-heuristic variants on the tuning bin (21--35 connectives), five seeds, mean $\pm$ std.}\label{tab:tune}
\resizebox{\columnwidth}{!}{\begin{tabular}{lccc}
\toprule
Method & solved $\uparrow$ & mean\_exp $\downarrow$ & $n$ \\
\midrule
Heuristic small & 0.88 $\pm$ 0.03 & 32.03 $\pm$ 5.17 & 5 \\
Heuristic large & \textbf{0.99 $\pm$ 0.01} & \textbf{24.85 $\pm$ 2.01} & 5 \\
Heuristic right & \underline{0.91 $\pm$ 0.01} & \underline{28.06 $\pm$ 3.03} & 5 \\
\bottomrule
\end{tabular}
}
\end{table}
